% Глава: 25-git
% Статус: наполнено (первый прогон).
\chapter{Git: системы контроля версий}
\label{ch:25-git}

\section{Зачем нужен контроль версий}

Версия кода --- это состояние проекта в какой-то
момент. Без контроля версий трудно: вернуть старый
вариант, понять, кто и что менял, работать над
несколькими задачами параллельно, делиться кодом.
\textbf{Git} --- распределённая система контроля
версий (VCS): полная история на каждой машине,
быстрые локальные операции, ветвление как базовый
механизм.

Работаем с Git на своей ВМ на \fpath{pglab.nsu.ru}: репозиторий
живёт на сервере, а удалённым хостом служит, например, GitHub
или внутренний Git на сети университета.

\section{Распределённая модель}

В централизованных системах (SVN) есть один центральный
сервер. В Git \emph{каждая копия} --- полноценное
хранилище со всей историей; ``сервер''
(\cmd{github.com}, \cmd{gitlab.com}) нужен для
координации, а не как обязательная точка. Работать
можно офлайн, коммиты локальные, а \cmd{push} ---
это синхронизация.

\section{Три области}

\begin{figure}[htb]
  \centering
  \begin{tikzpicture}[node distance=0.9cm]
    \node[data, minimum width=2.9cm] (wd) {working tree\\[0.2em] (files on disk)};
    \node[process, right=of wd, minimum width=2.9cm] (idx) {index\\[0.2em] (staging area)};
    \node[process, right=of idx, minimum width=2.9cm] (repo) {repository\\[0.2em] (commits)};
    \node[process, right=of repo, minimum width=2.9cm] (remote) {remote\\[0.2em] (origin)};
    \draw[flow] (wd) -- node[above, font=\small]{git add} (idx);
    \draw[flow] (idx) -- node[above, font=\small]{git commit} (repo);
    \draw[flow, <->] (repo) -- node[above, font=\small]{push / pull} (remote);
  \end{tikzpicture}
  \caption{Три области Git: рабочая директория, индекс, хранилище.}
  \label{fig:git-areas}
\end{figure}

\begin{itemize}
  \item \textbf{working tree} --- файлы на диске;
  \item \textbf{index} (staging area) --- набор
    изменений, которые попадут в следующий коммит;
  \item \textbf{repository} --- цепочка коммитов;
  \item \textbf{remote} --- удалённое хранилище
    (\cmd{origin}).
\end{itemize}

\section{Первые шаги}

\begin{minted}[bgcolor=gray!6]{bash}
git init                 # создать здесь новый репозиторий
git clone <url>          # скопировать существующий репозиторий
cd myrepo
\end{minted}

\section{Базовый цикл}

\begin{minted}[bgcolor=gray!6]{bash}
git status               # что изменилось
git add file.txt         # добавить один файл в индекс
git add .                # добавить все изменения этого каталога
git commit -m "message"  # создать коммит
git log --oneline        # краткая история
\end{minted}

\textbf{Коммит} --- атомарный снимок: кто, когда, что
изменил и почему. Текст коммита описывает \emph{что}
и \emph{зачем}.

\section{Ветвление и слияние}

\begin{minted}[bgcolor=gray!6]{bash}
git branch feature-x     # создать ветку
git switch feature-x     # переключиться на неё
git switch -c feature-y  # создать и переключиться за один шаг
git merge feature-x      # слить ветку в текущую
git branch -d feature-x  # удалить после слияния
\end{minted}

\textbf{merge} --- объединяет историю двух веток,
создавая коммит-слияние. \textbf{rebase} ---
``пересаживает'' вашу ветку на другую, история
линейнее, но коммиты переписываются.

\begin{minted}[bgcolor=gray!6]{bash}
git rebase main          # переиграть эту ветку поверх main
\end{minted}

\begin{tipbox}
Правило: \cmd{merge} для публичных веток (честная
история), \cmd{rebase} для личных веток до
публикации (чистая история). Никогда не
переписывайте уже опубликованную историю.
\end{tipbox}

\section{Работа с удалённым хранилищем}

\begin{minted}[bgcolor=gray!6]{bash}
git remote add origin <url>   # указать удалённый репозиторий
git push -u origin main       # отправить и задать upstream
git pull                      # забрать и слить с upstream
git fetch                     # только скачать, без слияния
git branch -r                 # список удалённых веток
\end{minted}

\section{\cmd{.gitignore}}

Файл \fpath{.gitignore} --- какие файлы Git
\emph{не} отслеживает:
\begin{minted}[bgcolor=gray!6]{bash}
# артефакты сборки
*.o
*.a
out/
build/

# редактор и IDE
.vscode/
.idea/
*.swp

# секреты и локальная конфигурация
.env
config.local.yml
\end{minted}

\section{Просмотр изменений и история}

\begin{minted}[bgcolor=gray!6]{bash}
git diff                 # недобавленные изменения
git diff --staged        # добавленные в индекс
git diff HEAD            # все изменения относительно последнего коммита
git log --oneline --graph --all   # наглядная история
git show <commit>        # что изменил коммит
\end{minted}

\section{Откат и исправление}

\begin{minted}[bgcolor=gray!6]{bash}
git reset --soft HEAD~1    # отменить последний коммит, оставить изменения в индексе
git reset --mixed HEAD~1   # отменить последний коммит, не добавляя в индекс
git reset --hard HEAD~1    # отменить последний коммит и выбросить изменения, осторожно!
git revert <commit>        # создать новый коммит, отменяющий указанный
\end{minted}

\cmd{reset} двигает указатель ветки; \cmd{revert} ---
безопасная операция, создающая новый коммит с
обратным эффектом. Для публичной истории используйте
\cmd{revert}.

\section{\cmd{.gitconfig} и идентификация}

\begin{minted}[bgcolor=gray!6]{bash}
git config --global user.name "Jane Doe"
git config --global user.email "jane@example.org"
git config --global init.defaultBranch main
git config --global core.editor vim
\end{minted}

\section{Типичные ошибки}
\begin{itemize}
  \item Забыли \cmd{git add} перед \cmd{commit} ---
    коммит пустой;
  \item В коммите ``fixed bug'' без объяснения ---
    невозможно понять, \emph{какой} баг и \emph{почему}
    так;
  \item \cmd{git push --force} на общую ветку ---
    история переписывается, у коллег конфликты;
  \item Секреты (пароли, токены) в репозитории ---
    добавьте в \cmd{.gitignore} и \emph{измените
    секрет};
  \item \cmd{git merge} и \cmd{git rebase} в одном
    рабочем процессе --- путаница; выберите один
    подход.
\end{itemize}

\section{Связь с остальным курсом}

Git --- это то, как делятся проекты в команде
(\nameref{ch:27-scripting-ide}) и что
интегрируют в CI (\nameref{ch:26-ci}). В Docker
(\nameref{ch:19-docker-basics}) образы часто
тегируются номером коммита.

\section{Задачи для самопроверки}
\begin{enumerate}
  \item Что такое коммит и чем он отличается от
    снимка в резервном копировании?
  \item Чем \cmd{merge} отличается от \cmd{rebase}?
    Когда что использовать?
  \item Как создать ветку, переключиться в неё и
    слить обратно?
  \item Что делает \cmd{git reset --hard} и почему он
    опасен?
  \item Какие файлы обычно попадают в
    \cmd{.gitignore}?
  \item Чем \cmd{pull} отличается от \cmd{fetch}?
\end{enumerate}
